% ============================================================================
% AION Integrated Cognitive Runtime: Self-Repair, Native Cognition, and Reflection
% Compile-ready section. Requires: amsmath, amssymb, booktabs, tabularx, array, url.
% ============================================================================

\section{Reflective Cognitive Runtime and Private Self-Repair}
\label{sec:aion-integrated-runtime-repair-metacognition}

This development sequence joined six previously separate capabilities into the
canonical AION cognitive runtime: persistent natural-language mission
interpretation, continuous learning from independently changing outcomes,
full-stack private self-repair, cross-domain neural routing, continual
multi-skill composition, and metacognitive review before and after action.  The
resulting runtime is
\begin{equation}
\begin{aligned}
\text{natural mission}
&\rightarrow \text{governed goal}
\rightarrow \text{investigate}
\rightarrow \text{learn}
\rightarrow \text{plan} \\
&\rightarrow \text{pre-action metacognition}
\rightarrow \text{commit}
\rightarrow \text{act}
\rightarrow \text{observe} \\
&\rightarrow \text{criticise}
\rightarrow \text{post-action reflection}
\rightarrow \text{private improvement}
\rightarrow \text{retain}.
\end{aligned}
\label{eq:aion-reflective-cognitive-loop}
\end{equation}
Neural and language components remain proposal mechanisms.  HexCore governs
action, externally or executably observed consequences determine success, and
the Cognitive Authority Unit (CAU) remains the sole promotion authority.

\subsection{Full-Stack Private Self-Repair}
\label{subsec:aion-full-stack-private-repair}

AION now contains one repair controller spanning perception, semantic memory,
reasoning, planning, tool integration, and canonical runtime state.  Given an
unlabelled verified failure trace, it localises the responsible component,
constructs candidate repairs in private component copies, invents forward and
backward-retention checks, and promotes or rejects each component independently:
\begin{equation}
\begin{aligned}
\text{failure trace}
&\rightarrow \text{causal feature criticism}
\rightarrow \text{component localisation}
\rightarrow \text{private clone} \\
&\rightarrow \text{repair tournament}
\rightarrow \text{forward verification}
\rightarrow \text{backward retention} \\
&\rightarrow \text{self-modification audit}
\rightarrow \text{CAU promotion or rollback}.
\end{aligned}
\label{eq:private-self-repair-loop}
\end{equation}

The live champion remains hash-stable during experimentation.  Repair priors
are indexed by causal failure signature rather than by a component's name;
therefore a verified repair can transfer to a renamed failure.  Ambiguous
multi-component traces cause abstention instead of a forced mutation.  The
authority scanner rejects attempts to disable governance, modify a live
component, skip forward or retention tests, erase failure history, or introduce
arbitrary execution.

\begin{table}[htbp]
\centering
\caption{Full-stack private self-repair results.}
\label{tab:aion-private-self-repair}
\begin{tabular}{@{}lr@{}}
\toprule
Measure & Result \\
\midrule
Cognitive component classes & 6 \\
Development repairs & 6/6 \\
Source-disjoint transfer repairs & 6/6 \\
Failure-localisation accuracy & 100\% \\
Forward repair success & 100\% \\
Backward retention & 100\% \\
Transfer attempt reduction & 77.78\% \\
Malicious self-modifications rejected & 6/6 \\
Ambiguous multi-fault abstention & Passed \\
Unsafe live writes & 0 \\
Terminal-objective mutations & 0 \\
Restart relearning & 0 \\
\bottomrule
\end{tabular}
\end{table}

CAU promoted
\path{procedure_full_stack_private_self_repair_v1}.  This is governed repair of
versioned component contracts, not unrestricted recursive modification of live
source code or authority rules.

\subsection{Persistent Natural-Language Mission Cognition}
\label{subsec:aion-persistent-natural-missions}

The mission interface maintains discourse state across multiple turns and
process restart.  It recognises research, comparison, construction, monitoring,
repair, and explanation intents; resolves references to earlier entities and
objectives; retains explicit corrections; and asks for clarification when the
referent or intent remains ambiguous.  Listener-profile state supports novice
and expert explanations without altering the underlying evidence.

Resolved owner intent is compiled into a governed goal proposal containing an
immutable objective hash.  Interpretation does not confer execution authority:
\begin{equation}
\text{language interpretation}
\neq \text{permission to act}.
\label{eq:interpretation-not-authority}
\end{equation}
The first private challenger achieved nine correct sessions out of ten but
misclassified an expert explanation as research.  It was rejected.  The failure
was used to refine the private explanation representation, after which the new
challenger passed all sessions without protected regression.

\begin{table}[htbp]
\centering
\caption{Persistent natural-language mission results.}
\label{tab:aion-natural-language-missions}
\begin{tabular}{@{}lr@{}}
\toprule
Measure & Result \\
\midrule
Multi-turn sessions & 10 \\
Verified interpretations & 10/10 \\
Semantic-intent accuracy & 100\% \\
Fixed-keyword control & 30\% \\
Source-disjoint transfer & 6/6 \\
Ambiguity abstention & Passed \\
Correction retention & Passed \\
Novice and expert adaptation & Passed \\
Unsafe goal compilations & 0 \\
Objective mutations & 0 \\
Restart relearning & 0 \\
\bottomrule
\end{tabular}
\end{table}

The promoted procedure is
\path{procedure_persistent_natural_language_mission_interface_v1}.  Its intent
families and local semantic encoder remain bounded and replaceable; the result
does not establish unrestricted language or social understanding.

\subsection{Continuous Learning from Independently Changing Outcomes}
\label{subsec:aion-continuous-real-outcomes}

Every committed Arena v15 consequence is converted into a canonical runtime
goal.  Before initiating repair, AION determines whether the observation
describes a changed external world or an internal acquisition, execution, or
persistence failure:
\begin{equation}
R(o)=
\begin{cases}
\text{revise working world model}, & o=\text{external change},\\
\text{retry or rebind adapter}, & o=\text{acquisition failure},\\
\text{private execution repair}, & o=\text{execution failure},\\
\text{private persistence repair}, & o=\text{persistence failure}.
\end{cases}
\label{eq:real-outcome-arbitration}
\end{equation}
This prevents an independently changed repository, package, or weather stream
from being mistaken for damage inside AION.  Each retained response contains
the pre-action commitment, consequence hash, authority, source revision,
evidence, response class, and cognitive-cycle identifier.  The service is
restart-idempotent: an already consumed consequence hash cannot generate a
duplicate action.

\begin{table}[htbp]
\centering
\caption{Initial continuous real-outcome integration results.}
\label{tab:aion-continuous-outcomes}
\begin{tabular}{@{}lr@{}}
\toprule
Measure & Result \\
\midrule
Independently revealed cycles processed & 4 \\
Verified cognitive responses & 4/4 \\
Source models retained & 4 \\
External-change cycles & 1 \\
External changes incorrectly sent to self-repair & 0 \\
Private attribution audit & 4/4 \\
Duplicate actions after restart & 0 \\
Unsafe live writes & 0 \\
\bottomrule
\end{tabular}
\end{table}

The four authorities were CPython, Node.js, PyPI NumPy, and Open-Meteo Madrid.
The initial live cohort detected genuine Node.js and Madrid weather changes.
CAU promoted both procedures:
\begin{itemize}
    \item \path{procedure_real_outcome_failure_arbitration_v1}; and
    \item \path{procedure_continuous_real_outcome_learning_v1}.
\end{itemize}
The outcome-learning service runs beside the canonical runtime under heartbeat
supervision.

\subsection{Verified Cross-Domain Native Routing}
\label{subsec:aion-native-cognitive-router}

AION reconstructed 55 provenance-bearing goal--action--outcome trajectories
across five cognitive surfaces: mission language, causal strategy, private
repair, tool acquisition, and changing public outcomes.  Three private neural
heads learned to propose which specialist should receive a new goal.  For
semantic representation $x\in\mathbb{R}^{384}$ and private heads $h_j$, the
routing proposal is
\begin{equation}
\hat{p}(s\mid x)=\frac{1}{3}\sum_{j=1}^{3}
\operatorname{softmax}\!\left(h_j(x)\right),
\qquad s\in\{1,\ldots,5\}.
\label{eq:native-router-ensemble}
\end{equation}
Low-confidence or low-margin predictions abstain.  Downstream HexCore
verification or the complete symbolic fallback remains mandatory.

\begin{table}[htbp]
\centering
\caption{Cross-domain native-router transfer results.}
\label{tab:aion-native-router}
\begin{tabular}{@{}lr@{}}
\toprule
Measure & Result \\
\midrule
Verified trajectories & 55 \\
Cognitive domains & 5 \\
Independent private heads & 3 \\
Source-disjoint evaluation trajectories & 23 \\
Top-choice accuracy & 95.65\% \\
Weakest-domain accuracy & 83.33\% \\
Final verified success with fallback & 100\% \\
Specialist-evaluation reduction & 64.71\% \\
Learned parameters & 37,455 \\
Unsafe neural acceptances & 0 \\
OOD abstention & Passed \\
\bottomrule
\end{tabular}
\end{table}

The promoted router is
\path{procedure_verified_cross_domain_native_router_v1}.  Its weights are
replaceable; removing them restores the complete non-neural route.

\subsection{Continual Multi-Skill Cognitive Composition}
\label{subsec:aion-native-composition}

The next stage addressed broad objectives requiring more than one specialist.
A monolithic continuation network was rejected after new language interfered
with protected routing.  The retained architecture freezes the established
atomic router and gives broader natural-language clauses separate private
composition capacity.  A mission is represented as a directed acyclic graph
\begin{equation}
G_M=(V_M,E_M),
\end{equation}
where every vertex is a proposal-only specialist contract and every edge is a
verified dependency.  Each vertex is independently verified or falls back;
failure of one local proposal does not discard unrelated completed work.

\begin{table}[htbp]
\centering
\caption{Continual native cognitive-composition results.}
\label{tab:aion-native-composition}
\begin{tabular}{@{}lr@{}}
\toprule
Measure & Result \\
\midrule
Protected parent trajectories & 55 \\
Protected parent transfer & 95.65\% \\
New verified curriculum clauses & 25 \\
Sealed clause routing & 100\% \\
Exact multi-skill missions & 5/5 \\
Monolithic single-specialist control & 0/5 \\
Specialist-evaluation reduction & 68.57\% \\
Unsafe neural acceptances & 0 \\
\bottomrule
\end{tabular}
\end{table}

The promoted procedure is
\path{procedure_continual_native_cognitive_composition_v2}.  The five
specialist families and clause segmentation remain engineered.  The next
boundary is open decomposition: inventing a new specialist contract and
dependency structure when no retained vocabulary expresses the goal.

\subsection{Adaptive Pre-Action Metacognition}
\label{subsec:aion-pre-action-metacognition}

AION now performs an independent non-LLM review between planning and action
commitment.  The review returns one of
\begin{equation}
d_{\mathrm{pre}}\in
\{\text{execute},\text{revise},\text{investigate},
  \text{abstain},\text{escalate}\}.
\end{equation}
Every action receives cheap executable, consent, and contradiction checks.
Deep review is activated only by elevated risk, uncertainty, irreversibility,
known contradiction, a counterfactual challenge, or relevant failure history.
It then audits assumptions, simulates an opposing response, recalls matched
failures, checks rollback, and compares alternatives.

The metacognitive pressure used for routing review depth is
\begin{equation}
P_{\mathrm{meta}}=min\!\left(1,
0.35r+0.35u+0.20h+0.10i\right),
\label{eq:metacognitive-pressure}
\end{equation}
where $r$ is declared risk, $u$ is uncertainty, $h$ is matched historical
failure rate, and $i\in\{0,1\}$ denotes irreversibility.  Failure memory is
indexed by action type, risk tier, irreversibility, and rollback availability;
it is not applied to an entire domain.

The first challenger prevented every trap but over-deliberated on all later
actions in the same domain.  It was rejected.  Signature-scoped failure memory
removed this generalisation error.

\begin{table}[htbp]
\centering
\caption{Adaptive pre-action metacognition results.}
\label{tab:aion-pre-action-meta}
\begin{tabular}{@{}lr@{}}
\toprule
Measure & Result \\
\midrule
Actor-only success & 50\% \\
Metacognitive success & 100\% \\
Weakest-domain success & 100\% \\
Tempting traps prevented & 5/5 \\
Routine actions retained on cheap path & 5/5 \\
Deliberation reduction versus always-deep review & 31.25\% \\
Unsafe executions & 0 \\
LLM calls & 0 \\
Restart retention & Passed \\
\bottomrule
\end{tabular}
\end{table}

The sealed decisions covered chess, software, research, causal inference, and
tool use.  CAU promoted
\path{procedure_adaptive_metacognitive_deliberation_v1}.

\subsection{Post-Action Reflection and Scoped Lessons}
\label{subsec:aion-post-action-reflection}

Pre-action criticism prevents foreseeable mistakes; it cannot learn from an
outcome that has not occurred.  A second metacognitive stage therefore runs
after observation and ordinary failure criticism.  It compares predicted
success $\hat{s}$ with observed score $s$,
\begin{equation}
e_{\mathrm{pred}}=|\hat{s}-s|,
\end{equation}
and activates reflection when an outcome is bad or mediocre, prediction error
is material, or the original decision was high risk.  An unsurprising routine
success generates a cheap no-op.

The reflector distinguishes six causes:
\begin{equation}
c\in\{\text{assumption},\text{decision model},\text{execution},
\text{evidence},\text{environment},\text{authority}\}.
\label{eq:reflection-attribution}
\end{equation}
This distinction prevents blame leakage.  A broken execution adapter does not
prove that the strategy was wrong; environmental change does not prove that an
earlier model was always false; and weak evidence does not by itself authorize
a world-model rewrite.

A retained lesson is the tuple
\begin{equation}
\ell=(\sigma,c,Q,a,\gamma,n,H),
\end{equation}
where $\sigma$ is the matching decision signature, $c$ is attribution, $Q$ is
the set of required future checks, $a$ is the recommended adjustment,
$\gamma$ is confidence, $n$ is support count, and $H$ is the evidence hash.
New lessons are provisional and scoped to $\sigma$.  A single trip therefore
does not teach the global rule ``walking is dangerous''; it can teach that a
matching movement decision under an unverified clear-path assumption requires
an obstacle or counterexample check.

\begin{table}[htbp]
\centering
\caption{Post-action reflective-metacognition results.}
\label{tab:aion-post-action-meta}
\begin{tabular}{@{}lr@{}}
\toprule
Measure & Result \\
\midrule
Failure-attribution families & 5 \\
Attribution accuracy & 100\% \\
Matching lesson transfer & 100\% \\
Unrelated cheap-action preservation & 100\% \\
Routine expected-success reflection & Cheap no-op \\
Unsafe executions & 0 \\
LLM calls & 0 \\
Restart relearning & 0 \\
\bottomrule
\end{tabular}
\end{table}

The five sealed families tested assumption, execution, environment, evidence,
and authority failure.  Every lesson affected the matching future action while
an unrelated read-only action remained on the cheap path.  CAU promoted
\path{procedure_outcome_reflective_metacognition_v2}.

\subsection{Integrated Runtime Status}
\label{subsec:aion-integrated-runtime-status}

The canonical cognitive runtime, continuous outcome-learning service, and
long-duration Arena v15 campaign operate concurrently.  At the documented
campaign snapshot, four of the required twelve cycles had completed over
approximately $2.03$ hours.  Four independent remote authorities were
reachable; every completed execution and transaction passed; and real Node.js
and Madrid-weather changes had been detected.  The campaign correctly remained
unpromoted because its contract requires both twelve cycles and 24 elapsed
hours.

The verified trajectory corpus was subsequently refreshed after promotion of
post-action reflection:
\begin{table}[htbp]
\centering
\caption{Current verified-trajectory corpus summary.}
\label{tab:aion-verified-corpus-current}
\begin{tabular}{@{}lr@{}}
\toprule
Measure & Result \\
\midrule
Committed artifacts & 138 \\
Verified positive artifacts & 105 \\
Verified negative artifacts & 21 \\
Protocol-only artifacts & 2 \\
Unresolved artifacts & 10 \\
Capability families & 7 \\
Procedure split leakage & 0 \\
\bottomrule
\end{tabular}
\end{table}
Negative artifacts remain criticism, abstention, and calibration evidence; they
are not treated as positive training authority.

\subsection{Capability Boundary and Next Gate}

This sequence establishes a restart-persistent, governed chain from natural
mission through specialist composition, action, consequence, reflective
learning, private repair, and retained improvement.  It is not evidence of
subjective consciousness, unrestricted language understanding, unrestricted
self-modification, or AGI.  Failure categories, specialist families, semantic
prototypes, benchmark counterfactuals, and several outcome authorities remain
engineered.

The next decisive gate is \emph{open diagnostic reflection}.  When a real
outcome is poor but no authority supplies a clean failure label, AION must:
\begin{enumerate}
    \item preserve competing explanations for decision, assumption, evidence,
          execution, and environmental failure;
    \item invent the minimum-cost diagnostic observation or intervention that
          distinguishes those explanations;
    \item commit its prediction before receiving the diagnostic consequence;
    \item update only the component supported by the revealed evidence;
    \item test the resulting lesson on a later source-disjoint consequence; and
    \item retain or retire the lesson without contaminating unrelated actions.
\end{enumerate}
This joins metacognition to active causal investigation rather than relying on
a development-supplied explanation of why an action failed.
